\clearpage
\fancyhead[L]{\small D'AGUILA.}
\hypertarget{native-546-middle}{}\label{native-546-middle}
\begin{center}\small Göttingische gelehrte Anzeigen. Stück 102. Seite 1009..1010. 1812 Juni 27.\\[3pt]\rule{.72\linewidth}{.3pt}\end{center}

\emph{Decouverte de l'orbite de la terre du point central de l'orbite du soleil, leur situation et leur forme; de la section du zodiaque, par le plan de l'équateur, et du mouvement concordant des deux globes. Avec figures. Par Mr.} \textsc{C. J. E. H. d'Aguila}, \emph{ancien élève du génie.} 1806. 432 Seiten in Octav. Paris. \emph{De l'imprimerie de Délance.}

Auf Verlangen holen wir hier eine kurze Anzeige dieses zwar schon vor 6 Jahren erschienenen, uns aber erst seit kurzem bekannt gewordenen, Buches nach. Wir haben hier einen neuen Anti-Copernicaner, der die Sonne jährlich einen Kreis um die Erde beschreiben und die Planeten sich um jene bewegen lässt (wiewol von diesen nur beiläufig und mit wenigen Worten die Rede ist). Der Erde hingegen ertheilt er ausser der Rotationsbewegung noch eine jährliche gleichförmige Bewegung in einem excentrischen Kreise in der Ebene des Aequators, und dieses Kreises Halbmesser setzt er $\frac{1}{59}$ von dem Halbmesser der jährlichen Sonnenbahn. Diese wichtige Entdeckung, welche der Verf. selbst unbegreiflich nennt und einer unmittelbaren Inspiration zuschreibt, daher er sie auch Gott selbst dedicirt, soll das neunzehnte Jahrhundert verherrlichen, und die Secten von \textsc{Copernicus}, \textsc{Kepler} und \textsc{Newton}, auf welche der Verf. mit tiefer Verachtung herabsieht, vernichten. Dies mag genug sein, um das Buch neugierigen Lesern, wenn auch nur als eine psychologische Merkwürdigkeit, zu empfehlen.

\clearpage
\fancyhead[L]{\small HARDINGS HIMMELSATLAS.}
\hypertarget{native-546-lower}{}\label{native-546-lower}
\begin{center}\small Göttingische gelehrte Anzeigen. Stück 158. Seite 1569..1571. 1812 October 3.\\[3pt]\rule{.72\linewidth}{.3pt}\end{center}

Die beiden ersten, in den Jahren 1809 und 1810 erschienenen, Lieferungen des grossen Himmelsatlasses unsers Hrn. Prof. \textsc{Harding} sind zu ihrer Zeit in unsern Blättern angezeigt; wir holen jetzt die Anzeige der dritten, im Jahre 1811 ausgegebenen, Lieferung nach, und verbinden damit zugleich die \SourceLeaf{547}Anzeige der vierten, welche so eben erschienen ist. Da wir über den Plan und das Eigenthümliche dieser Arbeit bereits bei der ersten Lieferung ausführlich berichtet haben, so brauchen wir hier nur zu versichern, dass die folgenden Lieferungen ganz in demselben Geiste und mit derselben Sorgfalt ausgearbeitet sind, und an Reichhaltigkeit und Nettigkeit des Stichs die erste Lieferung zum Theil noch übertreffen.

Die \emph{dritte} Lieferung besteht aus den Blättern VII, XVI, XVII und XVIII. Das Blatt Nr. VII erstreckt sich in gerader Aufsteigung von $239^\circ$ bis $281^\circ$, in der Abweichung von $2^\circ$ nordlich bis $32^\circ$ südlich, und enthält also Stücke von den Sternbildern Scorpion Schütze, Ophiuchus, Schlange und Sobieskisches Schild. Das Blatt Nr. XVI hat in der geraden Aufsteigung dieselbe Ausdehnung, wie Nr. VII, und geht in der Declination von $1^\circ$ südlich bis $33^\circ$nordlich; es enthält, ausser dem grössten Theile des Hercules, Stücke von der Schlange, dem Ophiuchus, dem Poniatovsky'schen Stier und der Leyer. Das ausgezeichnet reichhaltige Blatt Nr. XVII, zu welchem Hr. \textsc{von Lindenau} eine beträchtliche Anzahl Sternbestimmungen dem Verfasser geliefert hat, schliesst sich an das vorige, mit welchem es gleiche Ausdehnung nach der Declination hat, an, und geht bis $321^\circ$ gerader Aufsteigung. Es enthält, ausser dem ganzen Delphin, dem Pfeil und dem Fuchs mit der Gans, den grössten Theil des Füllen und des Adlers und Antinous, so wie Stücke vom Poniatovsky'schen Stier, der Schlange, dem Hercules, der Leyer, dem Schwan, Pegasus und Wassermann. Auf gleiche Weise schliesst sich wiederum das Blatt Nr. XVIII an das vorige und geht bis zur geraden Aufsteigung $1^\circ$. Wir finden darauf beinahe den ganzen Pegasus, nebst Stücken vom Schwan, Füllen, Wassermann, den Fischen und der Andromeda.

Die \emph{vierte Lieferung} gibt uns lauter südliche Zonen, wozu der Verf. bei seinem Aufenthalte auf der Mannheimer Sternwarte viele ergänzende Sterne selbst beobachtet hat, und besteht aus den Blättern XI, XII, XIII und XIV. Diese Blätter gehen alle von $1^\circ$ nordlich bis $33^\circ$ südlich, und schliessen sich so an einander, dass sie die Zone von $39^\circ$ bis $201^\circ$ in der geraden Aufsteigung umfassen. Wir haben also darauf, theilweise oder ganz, den Wallfisch, die Georgsharfe, das Laboratorium, den Eridanus, den Brandenburgischen Scepter, die Bildhauerwerkstatt, den Orion, Hasen, die Taube, den grossen Hund, das Schiff Argo, das Einhorn, die Buchdruckerwerkstatt, den Compass, die Luftpumpe, Hydra, den Sextanten, Löwen, die Jungfrau, den Becher, Raben und Centaur. Wir bemerken nur noch, dass von jetzt an die \textsc{Vandenhoeck}'sche Buchhandlung den Verschliess der Karten übernommen hat.
